\begin{frame}{전반부와 후반부, 무엇이 달라지나}
  \small
  \begin{columns}[T]
    \begin{column}{0.48\linewidth}
      \textbf{\color{ksablue}전반부 (중간고사 이전) --- 고전 ML}
      \begin{itemize}
        \item 모든 모델을 \kb{가설 $h$, 손실 $J$, 최적화} 세 질문으로 읽는 틀을 세운다.
        \item 특징(feature)은 \textbf{사람이 만든다}. 모델은 그 위에서 경계만 긋는다.
        \item 모델이 작아서 \textbf{손으로 계산하고 그림으로 확인}할 수 있다.
        \item 과적합, 편향-분산, train/val/test 분리처럼 \textbf{평가하는 법}을 익힌다.
      \end{itemize}
    \end{column}
    \begin{column}{0.48\linewidth}
      \textbf{\color{ksablue}후반부 (중간고사 이후) --- 신경망}
      \begin{itemize}
        \item 같은 세 질문이 그대로 쓰이지만, $h$가 \textbf{층을 쌓은 함수}가 된다.
        \item 특징을 \textbf{모델이 데이터에서 스스로 배운다} (표현학습).
        \item 이미지·시퀀스·그래프처럼 \textbf{데이터 구조에 맞춘 구조}를 설계한다.
        \item 맞히는 것을 넘어 \textbf{새 데이터를 만드는} 생성모델까지 간다.
      \end{itemize}
    \end{column}
  \end{columns}
  \vskip0.8em
  \kq{전반부에서 익힌 평가·최적화 도구가 후반부의 모든 신경망에서 그대로 다시 쓰인다.}
\end{frame}
